\documentclass[11pt]{article}
\usepackage[a4paper,margin=25mm]{geometry}
\usepackage{amsmath,amssymb,amsthm,booktabs,hyperref,xurl,microtype}
\setlength{\emergencystretch}{3em}
\hypersetup{hidelinks}
\newtheorem{theorem}{Theorem}
\newtheorem{lemma}[theorem]{Lemma}
\newtheorem{proposition}[theorem]{Proposition}
\newtheorem{corollary}[theorem]{Corollary}
\newtheorem{remark}[theorem]{Remark}
\newcommand{\Q}{\mathbb Q}
\newcommand{\Tr}{\operatorname{Tr}}
\newcommand{\Nm}{\operatorname{N}}
\title{A common quadratic norm polynomial\\
and the removal of private field coordinates}
\author{\(\mathbb Z\mathbb Q\) Diophantine Project}
\date{23 November 2026}
\begin{document}
\maketitle
\begin{abstract}
We give an exact compression of a family of quadratic norm quotients.
If all blocks have the same two field quantities \(K,Q\), then the
norms of \(K+RQ\), as the rational scalar \(R\) varies, are evaluations
of one quadratic polynomial. Three common rational coefficients replace
two private rational coordinates in every block. For twelve blocks this
removes 21 variables and 21 quadratic equations, including all zero-case
and occurrence checks. Applied to the completed scalar-chart tables,
the construction gives 1122 variables and 966 quadratic equations, or
1302 variables and 1122 equations. These counts precede the separate
nine-variable Pell reduction. The resulting quartics retain the
arithmetic interfaces of their predecessors and the unresolved source
universality assertion \(U_{\mathrm{em}}\); no unconditional
undecidability claim is made.
\end{abstract}

\section{The algebraic reduction}

Let \(L/\Q\) be a quadratic field, with its nontrivial automorphism
written \(z\mapsto\overline z\). Write \(\Tr\) and \(\Nm\) for
the field trace and norm. Suppose that \(K,Q\in L\) are common
quantities, that \(Q\ne0\), and that
\begin{equation}\label{traceQ}
                           \Tr(Q)\ne0.
\end{equation}
The letter \(K\) here denotes a field element, rather than a field.
For each \(1\le i\le k\), let \(R_i\in\Q^\times\) and
\(\tau_i\in\Q\), and suppose that a register \(R_{2,i}=R_i^2\)
is already present. The old private variable is a field element
\(J_i\in L\), constrained by
\begin{equation}\label{old}
 R_i(J_i-Q)=K,\qquad
 \tau_i\Tr(J_i-Q)=\Nm(J_i).
\end{equation}
The first equation has two rational components; the second has one.

\begin{lemma}[The common denominator is forced]\label{denom}
If \(k\ge1\) and \eqref{old} holds, then \(T:=\Tr(K)\ne0\).
\end{lemma}
\begin{proof}
If \(T=0\), the first equation and \(R_i\ne0\) give
\(\Tr(J_i-Q)=0\). The second gives \(\Nm(J_i)=0\), so
\(J_i=0\), because \(L\) is a field. The first equation is
then \(K=-R_iQ\), whence \(T=-R_i\Tr(Q)\ne0\), a contradiction.
\end{proof}

Introduce three common rational variables \(A,B,C\), and impose
\begin{align}
 A T&=\Nm(K),&
 B T&=\Tr(K\overline Q),&
 C T&=\Nm(Q).                                    \label{coefficients}
\end{align}
The new private equation is
\begin{equation}\label{new}
                  R_i\tau_i=A+B R_i+C R_{2,i}.
\end{equation}

\begin{theorem}[Exact norm compression]\label{compression}
Under the assumptions just stated, \eqref{old} for every \(i\)
and \eqref{coefficients}--\eqref{new} have exactly the same
projection to all retained coordinates. The deleted private values
are uniquely recovered by
\begin{equation}\label{recover}
                           J_i=Q+K/R_i.
\end{equation}
If \(K,Q\) are represented by their two rational components and
\(R_{2,i}\) is an existing coordinate, all new equations have total
degree at most two. The transformation removes \(2k-3\) rational
coordinates and \(2k-3\) rational quadratic equations.
\end{theorem}
\begin{proof}
For any rational \(R\), direct multiplication gives
\begin{equation}\label{normexpansion}
 \Nm(K+RQ)=\Nm(K)+R\Tr(K\overline Q)+R^2\Nm(Q).
\end{equation}
Suppose first that the old equations hold. Lemma~\ref{denom}
allows the unique definitions in \eqref{coefficients}. By
\eqref{recover},
\[
 \Tr(J_i-Q)=T/R_i,\qquad
 \Nm(J_i)=R_i^{-2}\Nm(K+R_iQ).
\]
Multiplying the old trace equation by \(R_i^2\), and using
\eqref{normexpansion}, yields
\[
 R_i\tau_i T=\Nm(K)+R_i\Tr(K\overline Q)+R_i^2\Nm(Q).
\]
Division by \(T\), together with \(R_{2,i}=R_i^2\), gives
\eqref{new}.

Conversely suppose that all the new equations hold. Their last
common equation gives \(CT=\Nm(Q)\ne0\), so \(T\ne0\)
without any additional inverse coordinate. Define \(J_i\) by
\eqref{recover}; this uses only the already required \(R_i\ne0\).
The first old equation follows immediately. Multiplying
\eqref{new} by \(T\) and applying the three common equations
gives the displayed norm equality, and division by \(R_i^2\)
recovers the second old equation. Thus every old private equation
is restored.

Norms and traces of quadratic-field elements have degree at most
two in their rational components. Hence all three common equations
are quadratic, as is \eqref{new}. We delete the \(2k\) rational
components of the \(J_i\), and add three rational coordinates.
The \(3k\) old rational residuals become \(k+3\), giving the
same saving \(2k-3\) in their number.
\end{proof}

\begin{remark}[The trace hypothesis has a precise role]
Condition \eqref{traceQ} is needed for the old-to-new direction.
For example, in \(L=\Q(\sqrt2)\), let \(k=1\), \(R_1=1\),
\(Q=\sqrt2\), \(K=-\sqrt2\), and \(J_1=0\). Then
\eqref{old} holds for every \(\tau_1\), but \(T=0\) and
\(CT=\Nm(Q)=-2\) is impossible. In the application below,
\(Q\) is a nonzero square in a real quadratic field, so
\eqref{traceQ} is automatic.
\end{remark}

\section{Application to the combined quadratic systems}

In the predecessor construction \cite{rechart}, the field is
\(L=\Q(\sqrt2)\), and the common quantities are
\[
 Q=f^2,\qquad f=1/x(H)\ne0,\qquad K=Z^2.
\]
The old field variable is \(J_i=\widehat J_i\), and
\(\tau_i=\widehat\tau_i\). The literal predecessor equations
are precisely \eqref{old}. Their other scalar and native-pole
equations use \(\widehat\tau_i\), the existing \(R_i^2\), and
the other native-pole variables, and do not use \(\widehat J_i\).
The preceding scalar-chart reduction removed all occurrences of
\(\widehat J_i\) from the common point relations.

Since \(f=f_0+f_1\sqrt2\ne0\),
\[
                   \Tr(Q)=2(f_0^2+2f_1^2)>0,
 \qquad \Nm(Q)=\Nm(f)^2>0.
\]
The reference chart already forces \(f\ne0\). The scalar
normalizer is the unchanged, strictly positive source rational
function of the distinguished input; thus \(R_i\ne0\) is
retained. Neither fact depends on the deleted private norm rows.
Consequently Theorem~\ref{compression} applies to the entire
rational zero set under the predecessor's stated interpretation.

\subsection{Every boundary is accounted for}

No inverse of \(A\), \(B\), \(\tau_i\), or a component of
\(K\) or \(Q\) is introduced. In particular \(B=0\) is
allowed. The only divisions in the proof are by \(R_i\) and
\(T\), whose nonvanishing has been proved. The new common
equation \(CT=\Nm(Q)\) also excludes \(Z=0\), since
\(K=Z^2=0\) would give \(T=0\). Conversely the old equations
exclude this case by Lemma~\ref{denom}. Thus the reduction
does not discard an old zero branch. All point charts, including
their finite tangent and alternate two-torsion fibers, and all
native-pole equations remain literally unchanged.

In this application \(T=\Tr(Z^2)>0\). It follows that
\(A>0\) and \(C>0\), but these inequalities are consequences,
not additional constraints or positivity wrappers.

\subsection{The exact rational equations and storage conventions}

Write
\[
 K=K_0+K_1\sqrt2,\qquad Q=Q_0+Q_1\sqrt2.
\]
The three common residuals, without an extra trace register, are
\begin{align}
 2K_0A-K_0^2+2K_1^2&=0,\
 2K_0B-2K_0Q_0+4K_1Q_1&=0,\
 2K_0C-Q_0^2+2Q_1^2&=0.                        \label{componentrows}
\end{align}
For each block the fourth kind of row is
\[
             R_i\widehat\tau_i-A-BR_i-C R_{2,i}=0.
\]
The existing common source field \(Q\) is stored as
\[
           Q_0=Y_{64},\qquad Q_1=2Y_{65},
\]
where \(Y_j\) denotes the stored value with physical source
label \(j\), rather than its final renumbered column.
The factor two in the imaginary component is essential.
Both interface maps give source labels 64 and 65 contact weight
one and scalar-normalizer weight zero; their physical-function
normalization cancels the common contact gauge. Accordingly the
displayed \(Y_{64},2Y_{65}\) are the components of the actual
field \(f^2\), with no additional factor of the contact gauge.
The literal retained square equations are
\[
 Y_{15}Y_{16}=Y_{65},\qquad
 Y_{15}^2+2Y_{16}^2=Y_{64}.
\]

The fields \(Z\) and \(\widehat J_i\), and the scalar
\(\widehat\tau_i\), have zero contact, normalizer, and parameter
weights. Their physical source labels are respectively
102--103, 109--110, and 111. The latter three labels are final
local columns 86--88 in both interfaces. The final local columns
89 and 90 are the unchanged \(R_i^2\) and \(R_i\) registers.
The scalar-chart predecessor has already made the two components
of \(K=Z^2\) common, with explicit quadratic square equations.

\section{Counts, scope, and reproducibility}

There are twelve integer-certificate blocks: eleven ordinary inputs
and one reciprocal input. Therefore \(2k-3=21\).
\begin{center}
\begin{tabular}{@{}lrrrr@{}}
\toprule
Interface & Old variables & Old rows & New variables & New rows\\
\midrule
General 163-coordinate interface &1143&987&1122&966\\
Specified 179-coordinate interface&1323&1143&1302&1122\\
\bottomrule
\end{tabular}
\end{center}
The complete new tables are stored in
\nolinkurl{undecidability/coefficients/jones1980_scalar_norm/}.
Their sums of squares have degree exactly four. All twelve input
certificates and the outer source are retained. In particular,
this compression introduces no extra ordinary rational \(abc\)
assumption and changes no height threshold or prime-selection
condition. It is an algebraic equivalence after the preceding
common-point construction.

The independent Pell reduction \cite{pell} removes nine additional
outer coordinates and three additional rows. Its affected rows and
variables are disjoint from those of the present transformation.
Their combination therefore gives 1113 variables with 963 rows,
or 1293 variables with 1119 rows. These combined counts refer to
the composition of the two certified transformations; the tables
in the directory above deliberately precede the Pell transformation.

The source universality assertion \(U_{\mathrm{em}}\) remains
unresolved, including the discrepancies in the consulted Jones
source displays. None of the algebraic equivalences proved here
establishes that assertion. All numerical quartic undecidability
claims must retain it and the predecessor arithmetic interfaces.

The deterministic compiler is
\nolinkurl{undecidability/code/build_scalar_norm_compression.py}.
The independent verifier
\nolinkurl{undecidability/verify_scalar_norm_compression.py}
does not import that compiler. It checks the norm identity, the
three old private rows in each block, the exact new rows, every
unchanged residual, the complete occurrence list of the 24
removed coordinates, all relevant source scales, the retained
square and normalizer rows, integrality of every coefficient,
quadratic degree, support, counts, and file hashes. Its finite
checks establish those algebraic and coefficient statements;
the automatic trace guard and rational-solution equivalence are
proved above.

\section{A broader norm-polynomial principle}

The reason for the saving is not specific to the numerical
elliptic curve. For a degree-\(d\) number field and common field
elements \(K,Q\), the expression
\[
                      \Nm(K+RQ)=\sum_{j=0}^d c_jR^j
\]
has \(d+1\) common rational coefficients. If a system contains
private field variables \(J_i=Q+K/R_i\), the norm-trace equation
\(\tau_i\Tr(J_i-Q)=\Nm(J_i)\) is equivalent, whenever
\(T=\Tr(K)\ne0\), to
\[
                  \tau_i R_i^{d-1}
                        =\sum_{j=0}^d (c_j/T)R_i^j.
\]
Thus the private field elements can be eliminated in favor of
one common list of scalar coefficients. If all needed powers
of \(R_i\) are already registered, these last private equations
are quadratic. In higher degree, the common coefficient
definitions themselves may require additional multiplication
registers; they must be counted before claiming a numerical
saving. In degree two there are exactly three coefficient
definitions and each is already quadratic, which is why the
present saving is exactly \(2k-3\) with no omitted register.

\enlargethispage{\baselineskip}
\begin{thebibliography}{9}\small
\bibitem{rechart}
\(\mathbb Z\mathbb Q\) Diophantine Project,
\emph{Factoring the scalar normalizer out of common point charts},
23 November 2026.
\nolinkurl{undecidability/papers/scalar_rechart_quartic_checkpoint.tex}.

\bibitem{scalar}
\(\mathbb Z\mathbb Q\) Diophantine Project,
\emph{A scalar-point sharing construction in 163 and 179 variables},
v82, 20 November 2026.
\nolinkurl{papers/quartic/scalar_point163_179_20261120_v82.tex}.

\bibitem{pell}
\(\mathbb Z\mathbb Q\) Diophantine Project,
\emph{An independent Pell-index proof and quadratic optimization},
23 November 2026.
\nolinkurl{undecidability/papers/pell_index_optimization.tex}.
\end{thebibliography}
\end{document}
